\chapter{Was offen bleibt}

Die drei Asymmetrien sind gezeigt. Die Zeit hat eine ausgezeichnete Grenze, das Jetzt, und keine beliebigen; der Raum hat beliebige Grenzen und keine ausgezeichnete, und wo er eine hat, den Ort, das Hier, ist sie Zeit im Raum oder Leib im Raum. Die Zeit hat eine Richtung, weil sie die Faser der Wirkungen ist; der Raum hat keine, und wo er Richtungen zeigt, kommen sie vom Leib, der sich in ihm bewegt, bis hinab zur Zahl seiner Dimensionen, die der Sinn eines Wirkungsgesetzes ist. Die Zeit hat ihre Ordnung an der Kausalkette; der Raum hat seine Ordnung von der Zeit und sein Maß von der Uhr, und was die Uhr braucht, ist ein Ding, kein Raummaß. Der gemeinsame Grund ist, dass die Differenzen des Raumes Differenzen des Gegenstandes sind und die der Zeit Differenzen des Anerkennens. Die Physik, die vom Gegenstand handelt, behandelt die Zeit deshalb nach dem Raum, mit Recht und mit einem Verlust, der im Vorzeichen angezeigt wird und in ihm nicht enthalten ist.

Die Rangfrage ist damit beantwortet, sofern sie Grenze, Richtung und Maß betrifft. Die Zeit gibt, der Raum nimmt auf. Aus der Zeit erhält man die Ordnung des Raumes, aus dem Raum keine Ordnung der Zeit. Wo der Raum wahrhaft ist, ist er Zeit; wo die Zeit sich aufhält, ist sie Raum.

Was nicht beantwortet ist, ist die Frage nach dem Raum selbst. Alle Bestimmungen, die er bisher erhalten hat, sind negative oder geliehene: das Nebeneinander der Fasern, die leere Zerstreuung des Außersichseins, das Leiden, die Weite, das Eingeräumtsein, die Isotropie. Aber an drei Stellen hat er widerstanden, und diese Stellen sind nicht zufällig dieselben, an denen die Ableitung am weitesten getrieben wurde. Reichenbach, der die Raummessung auf die Zeitmessung zurückführt, braucht einen materiellen Messkörper, den die Lichtgeometrie nicht liefert. Kant, der die Zeit zur Form aller Erscheinungen macht, braucht ein Beharrliches im Raum, ohne das die Zeit nicht einmal als Größe vorgestellt werden könnte. Heidegger, der die Räumlichkeit auf die Zeitlichkeit zurückführen wollte, hat es zurückgenommen. Baumann, der die Zeit aus der Dauer des Ich versteht, lässt den Raum als eine Anschauung stehen, die uns bleibt, wenn wir die Dinge wegdenken, und die in keine Kategorientafel passt.

Und Fink, der die Bewegung als Räumen und Zeitigen zugleich denkt, lässt den Raum dem Ursprung nach gleichrangig.

Man kann diese Stellen so lesen, dass der Raum nur der Rest ist, der übrig bleibt, wenn die Ableitung noch nicht vollendet ist. Man kann sie auch so lesen, dass der Raum etwas ist, das aus der Zeit nicht abgeleitet werden kann, weil es nicht von ihrer Art ist: kein Modus, kein Werden, keine Faser, sondern das, worin Fasern nebeneinander sein können, ohne dass eine die andere wäre. Zwischen diesen beiden Lesarten wird hier nicht entschieden. Es wird nur festgehalten, dass die zweite nicht durch das widerlegt ist, was gezeigt wurde. Alles, was gezeigt wurde, betrifft die Grenze, die Richtung und das Maß des Raumes; es betrifft nicht das, was begrenzt, gerichtet und gemessen wird.

Offen bleibt auch das, was Reichenbach vertagt hat: die Einsinnigkeit der Zeit, das Fortschreiten, das \zitat{ein evidentes Erlebnis und zugleich eine Aussage über die Wirklichkeit} ist. Es ist nun zu sehen, warum die Physik es nicht lösen kann. Das Fortschreiten ist der Wechsel des Modus, das Künftige, das gegenwärtig wird und vergeht. Die Physik hat den Modus nicht, und sie kann ihn nicht haben, ohne aufzuhören, vom Gegenstand zu handeln. Dass die Zeit fortschreitet, ist deshalb keine physikalische Aussage und keine bloß subjektive; es ist die Aussage, dass es Gegenwart gibt, und Gegenwart ist das, was ein Gegenstand nicht ist und wovon jeder Gegenstand her ist. Chibas Satz sagt es einfach: Um die Zukunft zu erfahren, muss man nur warten. Warum das so ist, warum die Zukunft von selbst kommt, die Vergangenheit nie und das Dort nur, wenn man geht, ist mit den drei Asymmetrien nicht erklärt. Es ist das, was sie voraussetzen.

Und offen bleibt schließlich das Und. Fink hat gefragt, ob Zeit und Raum nebeneinander oder gleichzeitig sind. Die Antwort, die sich ergeben hat, lautet: weder noch. Sie sind zugleich in dem Bewegten, das räumt und zeitigt, und in dem Bleibenden, das einräumt und weilt. Das Bewegte und das Bleibende sind der Leib und die Dinge, und was diese sind, davon war hier nur insoweit die Rede, als sie Zeit und Raum auseinanderlegen. Die Welt, die Fink als das nennt, worin beide beisammen sind, ist mit ihren beiden Seiten nicht beschrieben. Sie ist das, worin die Beschreibung stattfindet.
